\subsection{SUBSTITUTIONS IN A THEORY}

Let \(\mathfrak{c}\) be a theory, let \(A_{1}, A_{2}, \ldots, A_{n}\) be its explicit axioms, let \(\pmb{x}\) be a letter and \(\pmb{T}\) a term of \(\mathfrak{c}\) . Let \((\pmb{T}|\pmb{x})\mathfrak{c}\) be the theory whose signs and schemes are the same as those of \(\mathfrak{c}\) , and whose explicit axioms are \((\pmb{T}|\pmb{x})\pmb{A}_{1}, (\pmb{T}|\pmb{x})\pmb{A}_{2}, \ldots, (\pmb{T}|\pmb{x})\pmb{A}_{n}\) .

C2. Let A be a theorem in a theory \(\mathfrak{G}\) , let T be a term of \(\mathfrak{G}\) , and let x be a letter. Then \((T|x)A\) is a theorem in the theory \((T|x)\mathfrak{G}\) .

Let \(R_1, R_2, \ldots, R_n\) be a proof in \(\mathfrak{G}\) in which \(A\) appears. Consider the sequence \((T|x)R_1, (T|x)R_2, \ldots, (T|x)R_n\) , which is a sequence of relations in \(\mathfrak{G}\) by reason of CF8 (§ 1, no. 4). We shall show that this sequence is a proof in the theory \((T|x)\mathfrak{G}\) ; this will establish the criterion. If \(R_k\) is an implicit axiom of \(\mathfrak{G}\) , then \((T|x)R_k\) is again an implicit axiom of \(\mathfrak{G}\) (no. 1) and therefore of \((T/x)\mathfrak{G}\) . If \(R_k\) is an explicit axiom of \(\mathfrak{G}\) , then \((T|x)R_k\) is an explicit axiom of \((T|x)\mathfrak{G}\) . Finally, if \(R_k\) is preceded by relations \(R_i\) and \(R_j\) , where \(R_j\) is \(R_i \Longrightarrow R_k\) , then \((T|x)R_k\) is preceded by \((T|x)R_i\) and \((T|x)R_j\) , and the latter is identical with \((T|x)R_i \Longrightarrow (T|x)R_k\) (criterion CS5).

C3. Let A be a theorem of a theory G, let T be a term of G, and let x be a letter which is not a constant of G. Then \((T|x)A\) is a theorem of G.

This follows immediately from C2, because x does not feature in the explicit axioms of \(\mathfrak{G}\). More particularly, if \(\mathfrak{G}\) contains no explicit axioms, or if the explicit axioms contain no letters, then the criterion C3 applies without restriction on the letter x.

